% Source: https://artofproblemsolving.com/community/c3h1697376_relatively_prime_numbers
% Stable identity: aops-topic-1697376-q01
Every product is built from the primes $2,3,5,7$. Two products are relatively prime when they share no prime factor.

\textbf{In a coprime pair at most one product is even,} so at least one of the two is odd. The odd products use only $1,3,5,7$: they are $1\cdot3\cdot5=15$, $1\cdot3\cdot7=21$, $1\cdot5\cdot7=35$ and $3\cdot5\cdot7=105$. Any two of these share a prime, so in each coprime pair exactly one product is odd, and the other is even.

\textbf{Partners of each odd product.}
\begin{itemize}
\item $105=3\cdot5\cdot7$: the partner may use only the prime $2$, so it is a product of three of $1,2,4$: only $8$. One pair.
\item $35=5\cdot7$: the partner uses only $2$ and $3$, so it is a product of three of $1,2,3,4,6$ that is even. The products are $6,8,12,12,18,24,24,36,48,72$, with distinct values $6,8,12,18,24,36,48,72$. Eight pairs.
\item $21=3\cdot7$: the partner is a product of three of $1,2,4,5$: $8,10,20,40$. Four pairs.
\item $15=3\cdot5$: the partner is a product of three of $1,2,4,7$: $8,14,28,56$. Four pairs.
\end{itemize}
The total is $1+8+4+4=\boxed{17}$.
